\PassOptionsToPackage{table}{xcolor}
\documentclass[10pt,aspectratio=169]{beamer}
\newcommand{\LectureNumber}{32}
\input{course-theme.tex}
\title{Testing practice and course synthesis}
\subtitle{CSC103 Programming Fundamentals / Lecture 32}
\author{Dr. Muhammad Farhan}
\institute{Associate Professor of Computer Science\\Department of Computer Science\\COMSATS University Islamabad\\Sahiwal Campus}
\date{Fall 2026}
\begin{document}
\begin{frame}\titlepage\end{frame}

\begin{frame}[fragile,t]{The problem for this lecture}
\begin{columns}[T,onlytextwidth]\begin{column}{0.60\textwidth}
Design tests for a marks-file summary: normal data, one mark, threshold 50, empty file, malformed token and out-of-range mark.
\end{column}\begin{column}{0.35\textwidth}\small
Use a pure mean method that rejects empty or out-of-range data.
\end{column}\end{columns}
\note{Introduce the target problem before explaining the tools. Revisit it in the guided solution later.\par Syllabus reading: Reference material: JUnit Jupiter. CLO-4.}
\end{frame}

\begin{frame}[fragile,t]{Learning objectives}
\begin{columns}[T,onlytextwidth]\begin{column}{0.60\textwidth}
\begin{itemize}
\item Organise parameterised and regression tests
\item Distinguish unit and integration checks
\item Test a program that reads and summarises marks
\end{itemize}
\end{column}\begin{column}{0.35\textwidth}\small
CLO-4

Parameterised tests\par Isolation and fixtures\par Testing an integrated program\par Coverage and final review
\end{column}\end{columns}
\note{Explain how each objective supports the opening problem. The final questions check these objectives.\par Syllabus reading: Reference material: JUnit Jupiter. CLO-4.}
\end{frame}

\begin{frame}[fragile,t]{Parameterised tests}
\begin{itemize}
\item A parameterised test applies one behaviour check to several cases.
\item @CsvSource supplies small literal datasets.
\item Each case should add a reason to test.
\end{itemize}
\vspace{1mm}\begin{block}{Example}
\begin{lstlisting}
@ParameterizedTest
@CsvSource({"2,5,5", "-4,-2,-2", "3,3,3"})
void maximumCases(int a, int b, int expected) {
  assertEquals(expected, Numbers.max(a,b));
}
\end{lstlisting}
\end{block}
\note{Required imports: org.junit.jupiter.params.ParameterizedTest and org.junit.jupiter.params.provider.CsvSource. The supplied junit-jupiter dependency includes parameterised-test support.\par Syllabus reading: Reference material: JUnit Jupiter. CLO-4.}
\end{frame}

\begin{frame}[fragile,t]{Isolation and fixtures}
\begin{itemize}
\item Tests should avoid shared mutable state.
\item Temporary directories isolate file tests.
\item Resource cleanup matters even when an assertion fails.
\end{itemize}
\vspace{1mm}\begin{block}{Example}
\begin{lstlisting}
@TempDir Path temp;
Path file = temp.resolve("marks.txt");
// Write a controlled fixture for this test.
\end{lstlisting}
\end{block}
\note{The temporary directory facility belongs to Jupiter. A file-processing test checks several cooperating components and is an integration check, even if it is small.\par Syllabus reading: Reference material: JUnit Jupiter. CLO-4.}
\end{frame}

\begin{frame}[fragile,t]{Testing an integrated program}
\begin{itemize}
\item Separate parsing, validation and calculation.
\item Test pure calculations with small unit tests.
\item Use file fixtures to check the whole reading-and-summary path.
\end{itemize}
\vspace{1mm}\begin{block}{Example}
\begin{lstlisting}
Fixture: "40 50 60 70"
Expected mean: 55.0
Expected passes: 3
Invalid fixture: "50 bad 60"
\end{lstlisting}
\end{block}
\note{State whether invalid data rejects the whole file. The companion project rejects malformed or out-of-range marks and rejects empty data for a mean.\par Syllabus reading: Reference material: JUnit Jupiter. CLO-4.}
\end{frame}

\begin{frame}[fragile,t]{Coverage and final review}
\begin{itemize}
\item Line coverage says which lines executed, not whether outcomes were checked well.
\item A small changed condition can reveal weak tests.
\item Review algorithms, control flow, arrays and error handling together.
\end{itemize}
\vspace{1mm}\begin{block}{Example}
\begin{lstlisting}
Mutation thought experiment:
Change mark >= 50 to mark > 50.
Which test must fail?
Answer: the case at exactly 50.
\end{lstlisting}
\end{block}
\note{This is a manual reasoning exercise, not a claim that a mutation-testing tool ran. Final exam scheduling and marks remain governed by the supplied syllabus.\par Syllabus reading: Reference material: JUnit Jupiter. CLO-4.}
\end{frame}

\begin{frame}[fragile,t]{A test matrix connects requirements to evidence}
\begin{itemize}
\item Each requirement should have an observable result and a corresponding test.
\item Boundary cases expose comparisons that ordinary values may never challenge.
\item Failure cases should check the intended exception or message.
\end{itemize}
\vspace{1mm}\begin{block}{Example}
\begin{lstlisting}
Requirement: 50 counts as a pass.
Test input: {50}
Expected passing count: 1
\end{lstlisting}
\end{block}
\note{Explain the mechanism using the displayed example. Each requirement should have an observable result and a corresponding test. Boundary cases expose comparisons that ordinary values may never challenge. Failure cases should check the intended exception or message.\par Syllabus reading: Reference material: JUnit Jupiter. CLO-4.}
\end{frame}

\begin{frame}[fragile,t]{A file fixture isolates an integration test}
\begin{itemize}
\item A fixture supplies controlled input data for a test.
\item A temporary directory prevents collisions with a user's real files.
\item The test then checks the result of file reading and calculation together.
\end{itemize}
\vspace{1mm}\begin{block}{Example}
\begin{lstlisting}
Write "40 50 60 70" to a temporary file.
Read it through the production method.
Assert mean 55.0 and passing count 3.
\end{lstlisting}
\end{block}
\note{Explain the mechanism using the displayed example. A fixture supplies controlled input data for a test. A temporary directory prevents collisions with a user's real files. The test then checks the result of file reading and calculation together.\par Syllabus reading: Reference material: JUnit Jupiter. CLO-4.}
\end{frame}


\begin{frame}[fragile,t]{Integration tests cross component boundaries}
\begin{center}\begin{tikzpicture}
\node[box] (n0) at (0.0,0) {Temporary file};
\node[box] (n1) at (3.45,0) {Parser};
\node[box] (n2) at (6.9,0) {Summary methods};
\node[alt] (n3) at (10.350000000000001,0) {Assertions};
\draw[arr] (n0) -- (n1);
\draw[arr] (n1) -- (n2);
\draw[arr] (n2) -- (n3);
\end{tikzpicture}\end{center}
\begin{block}{What to notice}An integration test checks that collaborating components agree on data and behaviour.\end{block}
\note{Trace each labelled connection. An integration test checks that collaborating components agree on data and behaviour.}
\end{frame}
\begin{frame}[fragile,t]{Key distinctions}
\small\renewcommand{\arraystretch}{1.5}
\rowcolors{2}{pale}{white}
\begin{tabularx}{\textwidth}{>{\raggedright\arraybackslash}X>{\raggedright\arraybackslash}X>{\raggedright\arraybackslash}X}
\rowcolor{navy}
\color{white}\bfseries Requirement & \color{white}\bfseries Test input & \color{white}\bfseries Expected evidence\\
Mean calculation & 40 50 60 70 & 55.0\\
One-record handling & 80 & Mean 80.0\\
Passing threshold & 50 & Pass count 1\\
Empty mean policy & Empty file & Exception for mean\\
Format validation & 50 bad 60 & Parsing-related exception\\
Range validation & 101 & Range exception\\
\end{tabularx}
\vspace{3mm}\footnotesize These distinctions determine which operation or control structure fits the problem.
\note{\par Syllabus reading: Reference material: JUnit Jupiter. CLO-4.}
\end{frame}

\begin{frame}[fragile,t]{First example: execution}
\begin{lstlisting}
int[] marks = {40, 50, 60, 70};
int sum = 0, passes = 0;
for (int mark : marks) {
  sum += mark;
  if (mark >= 50) passes++;
}
System.out.println((double) sum / marks.length);
System.out.println(passes);
\end{lstlisting}
\note{This is the main body from Java\_Examples/Lecture32.java. The source file includes the complete class. Predict the result before the next slide.\par Syllabus reading: Reference material: JUnit Jupiter. CLO-4.}
\end{frame}

\begin{frame}[fragile,t]{First example: result}
\begin{columns}[T,onlytextwidth]\begin{column}{0.60\textwidth}
55.0\par 3\par A boundary test at 50 distinguishes \textgreater{}= from \textgreater{}.
\end{column}\begin{column}{0.35\textwidth}\small
Testing an integrated program

Coverage and final review
\end{column}\end{columns}
\note{Expected output and interpretation: 55.0
3
A boundary test at 50 distinguishes \textgreater{}= from \textgreater{}.\par Syllabus reading: Reference material: JUnit Jupiter. CLO-4.}
\end{frame}

\begin{frame}[fragile,t]{Guided practice}
\begin{columns}[T,onlytextwidth]\begin{column}{0.60\textwidth}
Design tests for a marks-file summary: normal data, one mark, threshold 50, empty file, malformed token and out-of-range mark.
\end{column}\begin{column}{0.35\textwidth}\small
Attempt the solution before continuing.

Use the definitions and examples from this lecture.
\end{column}\end{columns}
\note{Give students 8 minutes to attempt the problem. The following slides provide a complete solution. Original solution guidance: Normal: 40 50 60 70 gives mean55 and passes3. One 80 gives mean80 and passes1. One 50 passes. Empty data rejects a mean. "bad" fails parsing. -1 or 101 fails validation.\par Syllabus reading: Reference material: JUnit Jupiter. CLO-4.}
\end{frame}

\begin{frame}[fragile,t]{Solution strategy}
\begin{columns}[T,onlytextwidth]\begin{column}{0.60\textwidth}
\begin{itemize}
\item Use a pure mean method that rejects empty or out-of-range data.
\item Use a pass-count method with the inclusive threshold 50.
\item Compare a normal dataset with a single mark exactly on the boundary.
\end{itemize}
\end{column}\begin{column}{0.35\textwidth}\small
The next program uses fixed inputs so its result can be reproduced.
\end{column}\end{columns}
\note{Discuss why each step is needed. State the input assumptions before executing the program.\par Syllabus reading: Reference material: JUnit Jupiter. CLO-4.}
\end{frame}

\begin{frame}[fragile,t]{Complete solution (1/2)}
\begin{lstlisting}
public class Solution32 {
  public static void main(String[] args) throws Exception {
    int[] marks = {40, 50, 60, 70};
    System.out.println(mean(marks));
    System.out.println(passes(marks));
    System.out.println(passes(new int[]{50}));
  }
  static double mean(int[] marks) {
    if (marks.length == 0) throw new IllegalArgumentException("empty");
    long sum = 0;
    for (int mark : marks) {
\end{lstlisting}
\note{Complete runnable file: Java\_Examples/Solution32.java. Inputs are fixed in the program.  \par Syllabus reading: Reference material: JUnit Jupiter. CLO-4.}
\end{frame}

\begin{frame}[fragile,t]{Complete solution (2/2)}
\begin{lstlisting}
      if (mark < 0 || mark > 100) throw new IllegalArgumentException("range");
      sum += mark;
    }
    return (double) sum / marks.length;
  }
  static int passes(int[] marks) {
    int count = 0;
    for (int mark : marks) if (mark >= 50) count++;
    return count;
  }
}
\end{lstlisting}
\note{Complete runnable file: Java\_Examples/Solution32.java. Inputs are fixed in the program.  \par Syllabus reading: Reference material: JUnit Jupiter. CLO-4.}
\end{frame}

\begin{frame}[fragile,t]{Execution trace}
\small\renewcommand{\arraystretch}{1.5}
\rowcolors{2}{pale}{white}
\begin{tabularx}{\textwidth}{>{\raggedright\arraybackslash}X>{\raggedright\arraybackslash}X>{\raggedright\arraybackslash}X>{\raggedright\arraybackslash}X}
\rowcolor{navy}
\color{white}\bfseries Test & \color{white}\bfseries Expected & \color{white}\bfseries If \textgreater{}= becomes \textgreater{} & \color{white}\bfseries Detects change?\\
Passes in \{40,50,60,70\} & 3 & 2 & Yes\\
Passes in \{80\} & 1 & 1 & No\\
Passes in \{50\} & 1 & 0 & Yes\\
\end{tabularx}
\vspace{3mm}\footnotesize Follow the changing state in execution order.
\note{Manually derived trace for the complete solution. Use a pure mean method that rejects empty or out-of-range data. Use a pass-count method with the inclusive threshold 50. Compare a normal dataset with a single mark exactly on the boundary.\par Syllabus reading: Reference material: JUnit Jupiter. CLO-4.}
\end{frame}

\begin{frame}[fragile,t]{Solution result}
\begin{columns}[T,onlytextwidth]\begin{column}{0.60\textwidth}
55.0\par 3\par 1
\end{column}\begin{column}{0.35\textwidth}\small
Why this result follows

Compare a normal dataset with a single mark exactly on the boundary.
\end{column}\end{columns}
\note{Expected result: 55.0
3
1\par Syllabus reading: Reference material: JUnit Jupiter. CLO-4.}
\end{frame}

\begin{frame}[fragile,t]{A common incorrect approach}
A suite executes every line, but never checks the calculated average. Is the calculation verified?
\vspace{1mm}\begin{block}{Example}
\begin{lstlisting}
No. Executed lines can still produce incorrect results. Add explicit assertions with independently calculated expected averages.
\end{lstlisting}
\end{block}
\note{Ask students to predict the failure before discussing the explanation. The right side gives the correction.\par Syllabus reading: Reference material: JUnit Jupiter. CLO-4.}
\end{frame}

\begin{frame}[fragile,t]{Boundary and failure checks}
\begin{columns}[T,onlytextwidth]\begin{column}{0.60\textwidth}
Test 49, 50 and 51 as separate one-mark datasets.
\end{column}\begin{column}{0.35\textwidth}\small
Normal: 40 50 60 70 gives mean55 and passes3. One 80 gives mean80 and passes1. One 50 passes. Empty data rejects a mean. "bad" fails parsing. -1 or 101 fails validation.
\end{column}\end{columns}
\note{Use the specification to derive each expected result. Exercise solution: Normal: 40 50 60 70 gives mean55 and passes3. One 80 gives mean80 and passes1. One 50 passes. Empty data rejects a mean. "bad" fails parsing. -1 or 101 fails validation.\par Syllabus reading: Reference material: JUnit Jupiter. CLO-4.}
\end{frame}

\begin{frame}[fragile,t]{A complete file integration test (1/2)}
\begin{lstlisting}
package edu.cui.csc103;

import java.nio.file.*;
import java.nio.charset.StandardCharsets;
import org.junit.jupiter.api.Test;
import org.junit.jupiter.api.io.TempDir;
import static org.junit.jupiter.api.Assertions.*;

class FileSummaryTest {
  @TempDir Path folder;

\end{lstlisting}
\note{The temporary directory isolates this test. The assertions check both parsing and the computed summary. Numbers and MarksFile are supplied production classes in the companion project.\par Syllabus reading: Reference material: JUnit Jupiter. CLO-4.}
\end{frame}

\begin{frame}[fragile,t]{A complete file integration test (2/2)}
\begin{lstlisting}
  @Test
  void summarisesAControlledFile() throws Exception {
    Path file = folder.resolve("marks.txt");
    Files.writeString(file, "40 50 60 70", StandardCharsets.UTF_8);
    int[] marks = MarksFile.read(file);
    assertArrayEquals(new int[]{40,50,60,70}, marks);
    assertEquals(55.0, Numbers.mean(marks), 1e-9);
    assertEquals(3, Numbers.passes(marks));
  }
}
\end{lstlisting}
\note{The temporary directory isolates this test. The assertions check both parsing and the computed summary. Numbers and MarksFile are supplied production classes in the companion project.\par Syllabus reading: Reference material: JUnit Jupiter. CLO-4.}
\end{frame}

\begin{frame}[fragile,t]{Check your understanding}
\begin{columns}[T,onlytextwidth]\begin{column}{0.60\textwidth}
Which tests should fail if \textgreater{}=50 becomes \textgreater{}50? What does a file integration test cover beyond a pure unit test?
\end{column}\begin{column}{0.35\textwidth}\small
Write a prediction and explain the rule that supports it.
\end{column}\end{columns}
\note{Answer: A case containing a mark exactly 50 with a pass-count assertion. The integration test also checks file opening, encoding, token parsing and cooperation between methods.\par Syllabus reading: Reference material: JUnit Jupiter. CLO-4.}
\end{frame}

\begin{frame}[fragile,t]{Answer and explanation}
\begin{columns}[T,onlytextwidth]\begin{column}{0.60\textwidth}
A case containing a mark exactly 50 with a pass-count assertion. The integration test also checks file opening, encoding, token parsing and cooperation between methods.
\end{column}\begin{column}{0.35\textwidth}\small
No. Executed lines can still produce incorrect results. Add explicit assertions with independently calculated expected averages.
\end{column}\end{columns}
\note{Reconnect the answer to the relevant definition rather than treating it as a fact to memorise.\par Syllabus reading: Reference material: JUnit Jupiter. CLO-4.}
\end{frame}


\begin{frame}[fragile,t]{Compare cases and predict outcomes}
\small\renewcommand{\arraystretch}{1.5}\rowcolors{2}{pale}{white}
\begin{tabularx}{\textwidth}{>{\raggedright\arraybackslash}X>{\raggedright\arraybackslash}X>{\raggedright\arraybackslash}X}\rowcolor{navy}
\color{white}\bfseries Fixture & \color{white}\bfseries Expected outcome & \color{white}\bfseries Boundary exercised\\
40 60 & Mean 50.0; passes 1 & File through summary\\
50 bad & Parsing failure & Malformed file token\\
Empty file & Empty array; mean rejected & No-data contract\\
\end{tabularx}\par\vspace{3mm}\begin{exampleblock}{Discuss}Explain each expected result before running the example. Which case would reveal a wrong assumption?\end{exampleblock}
\note{Read the first two columns and invite predictions before discussing the outcome.}
\end{frame}

\begin{frame}[fragile,t]{Apply the idea}
\begin{alertblock}{Independent practice}Plan an integration test for a temporary file containing 40 and 60. List the assertions that distinguish successful reading from successful calculation.\end{alertblock}\vspace{3mm}\begin{enumerate}\item Work through the input by hand.\item Write or correct the program.\item Justify the expected result.\end{enumerate}\note{Allow 3–5 minutes, then compare answers in pairs. Plan an integration test for a temporary file containing 40 and 60. List the assertions that distinguish successful reading from successful calculation.}
\end{frame}

\begin{frame}[fragile,t]{Solution and reasoning}
\begin{lstlisting}
Parsed array: assertArrayEquals(new int[]{40, 60}, marks)
Mean: assertEquals(50.0, Numbers.mean(marks), 1e-9)
Passes: assertEquals(1, Numbers.passes(marks))
\end{lstlisting}\vspace{1mm}\small\begin{itemize}
\item The array assertion checks parsing and record order.
\item The mean and pass assertions check the downstream computation.
\item Use an isolated temporary directory so tests do not depend on a personal file.
\end{itemize}\note{Answer: The array assertion checks parsing and record order. The mean and pass assertions check the downstream computation. Use an isolated temporary directory so tests do not depend on a personal file.}
\end{frame}

\begin{frame}[fragile,t]{Integration checks for student{-}record files}
\begin{itemize}
\item A temporary fixture isolates the test from personal files and previous runs.
\item Check parsing separately from an aggregate calculation.
\item Malformed fields and out{-}of{-}range scores are distinct failure cases.
\end{itemize}
\vspace{3mm}\footnotesize\textcolor{accent}{Source: supplied ISB campus slides. See speaker notes and ISB\_Source\_Map.md.}
\note{Adapted from the supplied ISB campus folder: File Handling.pptx, slides 14, 15, 18. Instructor{-}authored integration{-}test extension of the source records. List\textless{}String\textgreater{} is a supporting library use, not a new assessed data{-}structure unit.}
\end{frame}

\begin{frame}[fragile,t]{Worked problem: Integration checks for student{-}record files}
\begin{block}{Task}Write the two source records to a temporary UTF{-}8 file, parse the scores, and assert scores \{90,85\} and mean 87.5.\end{block}
\small Input: Values are fixed in the program for a reproducible demonstration.\par\vspace{3mm}\begin{exampleblock}{Before running}Predict the output and identify the variables that change.\end{exampleblock}
\note{Adapted from the supplied ISB campus folder: File Handling.pptx, slides 14, 15, 18. Instructor{-}authored integration{-}test extension of the source records. List\textless{}String\textgreater{} is a supporting library use, not a new assessed data{-}structure unit.}
\end{frame}

\begin{frame}[fragile,t]{Complete ISB example (1/2)}
\begin{lstlisting}
public class IsbLecture32 {
  public static void main(String[] args) throws Exception {
    java.nio.file.Path file = java.nio.file.Files.createTempFile("isb-test-", ".txt");
    java.nio.file.Files.writeString(file, "John T Smith 90\nEric K Jones 85\n",
      java.nio.charset.StandardCharsets.UTF_8);
    int[] scores = readScores(file);
    org.junit.jupiter.api.Assertions.assertArrayEquals(new int[]{90, 85}, scores);
    double mean = (scores[0] + scores[1]) / 2.0;
    org.junit.jupiter.api.Assertions.assertEquals(87.5, mean, 1e-9);
    System.out.println("Record integration passed");
  }
  static int[] readScores(java.nio.file.Path file) throws java.io.IOException {
    java.util.List<String> lines = java.nio.file.Files.readAllLines(
\end{lstlisting}
\note{Adapted from the supplied ISB campus folder: File Handling.pptx, slides 14, 15, 18. Instructor{-}authored integration{-}test extension of the source records. List\textless{}String\textgreater{} is a supporting library use, not a new assessed data{-}structure unit. Complete file: ISB\_Java\_Examples/IsbLecture32.java.}
\end{frame}

\begin{frame}[fragile,t]{Complete ISB example (2/2)}
\begin{lstlisting}
      file, java.nio.charset.StandardCharsets.UTF_8);
    int[] scores = new int[lines.size()];
    for (int i = 0; i < lines.size(); i++) {
      String[] fields = lines.get(i).trim().split("\\s+");
      if (fields.length != 4) throw new IllegalArgumentException("four fields required");
      int score = Integer.parseInt(fields[3]);
      if (score < 0 || score > 100) throw new IllegalArgumentException("score range");
      scores[i] = score;
    }
    return scores;
  }
}
\end{lstlisting}
\note{Adapted from the supplied ISB campus folder: File Handling.pptx, slides 14, 15, 18. Instructor{-}authored integration{-}test extension of the source records. List\textless{}String\textgreater{} is a supporting library use, not a new assessed data{-}structure unit. Complete file: ISB\_Java\_Examples/IsbLecture32.java.}
\end{frame}

\begin{frame}[fragile,t]{Worked trace and expected result}
\small\renewcommand{\arraystretch}{1.4}\rowcolors{2}{pale}{white}
\begin{tabularx}{\textwidth}{>{\raggedright\arraybackslash}X>{\raggedright\arraybackslash}X>{\raggedright\arraybackslash}X}\rowcolor{navy}
\color{white}\bfseries Fixture & \color{white}\bfseries Expected parsing & \color{white}\bfseries Expected next step\\
Two source records & [90, 85] & Mean 87.5\\
John T Smith bad & Reject non{-}integer & No aggregate\\
John T Smith 101 & Reject range & No aggregate\\
\end{tabularx}\par\vspace{2mm}
\begin{block}{Expected console output}\ttfamily\footnotesize Record integration passed\end{block}
\note{Adapted from the supplied ISB campus folder: File Handling.pptx, slides 14, 15, 18. Instructor{-}authored integration{-}test extension of the source records. List\textless{}String\textgreater{} is a supporting library use, not a new assessed data{-}structure unit. Trace and expected values independently worked for this edition.}
\end{frame}

\begin{frame}[fragile,t]{Extend the ISB example}
\begin{alertblock}{Practice}Why is checking only mean 87.5 weaker than also checking the parsed array?\end{alertblock}\vspace{3mm}\begin{exampleblock}{Explain your reasoning}Different incorrect arrays can share the same mean. The array assertion checks record count, values and order.\end{exampleblock}
\note{Adapted from the supplied ISB campus folder: File Handling.pptx, slides 14, 15, 18. Instructor{-}authored integration{-}test extension of the source records. List\textless{}String\textgreater{} is a supporting library use, not a new assessed data{-}structure unit. Answer: Different incorrect arrays can share the same mean. The array assertion checks record count, values and order.}
\end{frame}
\begin{frame}[fragile,t]{Lecture summary}
\begin{columns}[T,onlytextwidth]\begin{column}{0.60\textwidth}
\begin{itemize}
\item parameterised and regression tests
\item unit and integration checks
\item a program that reads and summarises marks
\end{itemize}
\end{column}\begin{column}{0.35\textwidth}\small
\begin{itemize}
\item Use a pure mean method that rejects empty or out-of-range data.
\item Use a pass-count method with the inclusive threshold 50.
\item Compare a normal dataset with a single mark exactly on the boundary.
\end{itemize}
\end{column}\end{columns}
\note{Ask students to give one concrete example for the concepts listed. Use the worked solution to connect the topics.\par Syllabus reading: Reference material: JUnit Jupiter. CLO-4.}
\end{frame}

\begin{frame}[fragile,t]{After-class practice}
\begin{columns}[T,onlytextwidth]\begin{column}{0.60\textwidth}
Prepare a small marks-analysis project with methods, arrays, file I/O and a documented test suite. Use the companion project as a starting point.
\end{column}\begin{column}{0.35\textwidth}\small
Syllabus reading\par Reference material: JUnit Jupiter

Next: course revision and final-exam preparation
\end{column}\end{columns}
\note{Reading labels follow the supplied outline. Check topic headings against the available textbook edition. Require a program and expected results for the practice task.\par Syllabus reading: Reference material: JUnit Jupiter. CLO-4.}
\end{frame}
\end{document}
